\frfilename{mt114.tex}
\versiondate{14.6.05}
\copyrightdate{1994}

\def\chaptername{Ruang ukur}
\def\sectionname{Ukuran Lebesgue pada $\Bbb R$}

\newsection{114}
\def\headlinesectionname{Ukuran Lebesgue pada ${\eightBbb R}$}

Sesudah gagasan-gagasan yang sangat abstrak dalam \S\S111-113, kita
sangat memerlukan sebuah contoh ruang ukur yang taktrivial. Contoh yang
jauh lebih penting daripada contoh-contoh lain adalah garis real dengan
ukuran Lebesgue, dan
sekarang saya akan menguraikan ukuran ini (114A-114E), beserta beberapa
sifat dasarnya.

Gagasan-gagasan utama dalam bagian ini diulangi dalam \S115, dan jika
Anda pernah menjumpai ukuran Lebesgue sebelumnya, atau merasa yakin
dapat menangani ruang berdimensi dua dan tiga sekaligus melakukan
analisis yang cukup sulit, Anda dapat langsung menuju bagian tersebut
dan kembali ke bagian ini hanya ketika diberikan rujukan khusus.

\leader{114A}{Definisi (a)} Untuk keperluan bagian ini, sebuah {\bf
interval setengah terbuka} dalam $\Bbb R$ adalah himpunan berbentuk
$\coint{a,b}=\{x:a\le x<b\}$, dengan $a$, $b\in\Bbb R$.

\cmmnt{Perhatikan bahwa saya mengizinkan $b\le a$ dalam rumus ini;
dalam hal ini $\coint{a,b}=\emptyset$ (lihat 1A1A).}

\header{114Ab}{\bf (b)}\cmmnt{ Jika $I\subseteq \Bbb R$ adalah interval
setengah terbuka, maka
$I=\emptyset$ atau $I=\coint{\inf I,\sup I}$, sehingga titik-titik
ujungnya terdefinisi dengan baik.} Kita\cmmnt{ karena itu dapat}
mendefinisikan {\bf panjang} $\lambda I$ dari interval setengah terbuka
$I$ dengan menetapkan

\Centerline{$\lambda\emptyset = 0$,
\quad $\lambda\coint{a,b}=b-a$ jika $a<b$.}

\leader{114B}{Lema} Jika $I\subseteq\Bbb R$ adalah interval setengah
terbuka dan $\langle I_j\rangle_{j\in\Bbb N}$ adalah barisan interval
setengah terbuka yang menutupi $I$, maka
$\lambda I\le\sum_{j=0}^{\infty}\lambda I_j$.

\proof{{\bf (a)}
Jika $I=\emptyset$, tentu saja
$\lambda I=0\le\sum_{j=0}^{\infty}\lambda I_j$. Jika tidak, ambil
$I=\coint{a,b}$, dengan $a<b$. Untuk setiap $x\in \Bbb R$, misalkan
$H_x$ adalah setengah garis $\ooint{-\infty,x}$, dan tinjau himpunan

\Centerline{$A=\{x:a\le x\le b,\,
x-a\le\sum_{j=0}\sp{\infty}\lambda(I_j\cap
H_x)\}$.}

\noindent (Perhatikan bahwa jika $I_j=\coint{c_j,d_j}$, maka $I_j\cap
H_x=\coint{c_j,\min(d_j,x)}$, sehingga $\lambda(I_j\cap H_x)$ selalu
terdefinisi.) Kita mempunyai $a\in A$ (karena
$a-a=0\le\sum_{j=0}\sp{\infty}\lambda(I_j\cap
H_a)$) dan tentu saja $A\subseteq[a,b]$, sehingga $c=\sup A$
terdefinisi dan termasuk dalam $[a,b]$.

\medskip

{\bf (b)} Sekarang kita dapati bahwa $c\in A$.

\prooflet{

$$\eqalign{\Prf\ c-a&=\sup_{x\in A}x-a\cr
&\le\sup_{x\in A}\sum_{j=0}^{\infty}\lambda(I_j\cap H_x)
\le\sum_{j=0}^{\infty}\lambda(I_j\cap H_c). \text{ \Qed}\cr}$$
}

\medskip

{\bf (c)} \Quer\ Andaikan, jika mungkin, bahwa $c<b$. Maka $c\in
\coint{a,b}$, sehingga terdapat suatu $k\in\Bbb N$ sedemikian sehingga
$c\in I_k$. Nyatakan $I_k$ sebagai $\coint{c_k,d_k}$; maka
$x=\min(d_k,b)>c$. Untuk setiap $j$,
$\lambda(I_j\cap H_x)\ge\lambda(I_j\cap H_c)$, sedangkan

\Centerline{$\lambda(I_k\cap H_x)=\lambda(I_k\cap H_c)+x-c$.}

\noindent Jadi

$$\eqalign{\sum_{j=0}^{\infty}\lambda(I_j\cap H_x)
&\ge\sum_{j=0}^{\infty}\lambda(I_j\cap H_c)+x-c\cr
&\ge c-a+x-c
=x-a,\cr}$$

\noindent sehingga $x\in A$; tetapi $x>c$ dan $c=\sup A$. \Bang

\medskip

{\bf (d)} Kita menyimpulkan bahwa $c=b$, sehingga $b\in A$ dan

\Centerline{$b-a\le\sum_{j=0}^{\infty}\lambda(I_j\cap
H_b)\le\sum_{j=0}^{\infty}\lambda I_j$,}

\noindent sebagaimana diklaim.
}%end of proof of 114B


\vleader{36pt}{114C}{Definisi} Sekarang\cmmnt{, dan untuk seluruh sisa
bagian ini,} definisikan $\theta:\Cal P\Bbb R\to[0,\infty]$ dengan
menetapkan

$$\eqalign{\theta A=\inf\{\sum_{j=0}\sp{\infty}\lambda I_j:\langle
I_j\rangle_{j\in\Bbb N}
\text{ adalah barisan }&\text{interval setengah terbuka}\cr
&\text{sedemikian sehingga }
A\subseteq\bigcup_{j\in\Bbb N}I_j\}.\cr}$$

\noindent\cmmnt{Perhatikan bahwa setiap $A$ dapat ditutupi oleh suatu
barisan interval setengah terbuka --- misalnya,
$A\subseteq\bigcup_{n\in\Bbb N}\coint{-n,n}$; sehingga jika kita
menafsirkan jumlah-jumlahnya dalam $[0,\infty]$, seperti dalam 112Bc di
atas, kita selalu mempunyai himpunan tak kosong yang infimumnya dapat
diambil, dan $\theta A$ selalu terdefinisi dalam $[0,\infty]$.
Fungsi ini} $\theta$ disebut {\bf ukuran luar Lebesgue} pada
$\Bbb R$\cmmnt{; frasa ini dibenarkan oleh (a) dari proposisi
berikutnya}.



\leader{114D}{Proposisi} (a) $\theta$ adalah ukuran luar pada $\Bbb R$.

(b) $\theta I=\lambda I$ untuk setiap interval setengah terbuka
$I\subseteq\Bbb R$.

\proof{{\bf (a)(i)} $\theta$ mengambil nilai dalam $[0,\infty]$
karena setiap $\theta A$ merupakan infimum suatu subhimpunan tak kosong
dari $[0,\infty]$.

\medskip

\quad{\bf (ii)} $\theta\emptyset=0$ karena (misalnya) jika kita
menetapkan $I_j=\emptyset$ untuk setiap $j$, maka setiap $I_j$ adalah
interval setengah terbuka (menurut kesepakatan yang saya gunakan) dan
$\emptyset\subseteq\bigcup_{j\in\Bbb N}I_j$,
$\sum_{j=0}\sp{\infty}\lambda I_j=0$.

\medskip

\quad{\bf (iii)} Jika $A\subseteq B$, maka setiap kali
$B\subseteq\bigcup_{j\in\Bbb N}I_j$ kita mempunyai
$A\subseteq\bigcup_{j\in\Bbb N}I_j$, sehingga $\theta A$ adalah
infimum suatu himpunan yang sekurang-kurangnya sama besar dengan
himpunan yang terlibat dalam definisi $\theta B$, dan
$\theta A\le\theta B$.

\medskip

\quad{\bf (iv)} Sekarang andaikan bahwa
$\langle A_n\rangle_{n\in\Bbb N}$ adalah barisan subhimpunan
$\Bbb R$, dengan gabungan $A$. Untuk setiap $\epsilon>0$, kita dapat
memilih, bagi setiap $n\in\Bbb N$, suatu barisan
$\langle I_{nj}\rangle_{j\in\Bbb N}$ yang terdiri atas interval
setengah terbuka sedemikian sehingga
$A_n\subseteq\bigcup_{j\in\Bbb N}I_{nj}$ dan
$\sum_{j=0}\sp{\infty}\lambda I_{nj}\le\theta A_n+2^{-n}\epsilon$.
(Anda mungkin perlu memeriksa bahwa perumusan ini sah, baik ketika
$\theta A_n$ berhingga maupun tak hingga.) Menurut 111F(b-ii), terdapat
bijeksi dari $\Bbb N$ ke $\Bbb N\times\Bbb N$; nyatakan bijeksi ini
dalam bentuk $m\mapsto(k_m,l_m)$. Maka
$\sequence{m}{I_{k_m,l_m}}$ adalah barisan interval setengah terbuka,
dan

\Centerline{$A\subseteq\bigcup_{m\in\Bbb N}I_{k_m,l_m}$.}

\noindent \Prf\ Jika $x\in A=\bigcup_{n\in\Bbb N}A_n$, harus ada
$n\in\Bbb N$ sedemikian sehingga
$x\in A_n\subseteq\bigcup_{j\in\Bbb N}I_{nj}$, sehingga terdapat
$j\in\Bbb N$ sedemikian sehingga $x\in I_{nj}$. Karena
$m\mapsto(k_m,l_m)$ surjektif, terdapat $m\in\Bbb N$ sedemikian sehingga
$k_m=n$ dan $l_m=j$; dalam hal ini $x\in I_{k_m,l_m}$.\ \Qed

Selanjutnya,

\Centerline{$\sum_{m=0}\sp{\infty}\lambda I_{k_m,l_m}
\le\sum_{n=0}\sp{\infty}\sum_{j=0}\sp{\infty}\lambda I_{nj}$.}

\noindent \Prf\ Jika $M\in\Bbb N$, maka
$N=\max(k_0,k_1,\ldots,k_M,l_0,l_1,\ldots,l_M)$ berhingga; karena
setiap $\lambda I_{nj}$ lebih besar daripada atau sama dengan $0$, dan
setiap pasangan $(n,j)$ dapat muncul paling banyak sekali sebagai suatu
$(k_m,l_m)$,

\Centerline{$\sum_{m=0}^M\lambda I_{k_m,l_m}
\le\sum_{n=0}^N\sum_{j=0}^N\lambda I_{nj}
\le\sum_{n=0}^N\sum_{j=0}^{\infty}\lambda I_{nj}
\le\sum_{n=0}^{\infty}\sum_{j=0}^{\infty}\lambda I_{nj}$.}

\noindent Jadi

\Centerline{$\sum_{m=0}^{\infty}\lambda I_{k_m,l_m}
=\lim_{M\to\infty}\sum_{m=0}^M\lambda I_{k_m,l_m}
\le\sum_{n=0}^{\infty}\sum_{j=0}^{\infty}\lambda I_{nj}$. \Qed}

Dengan demikian

$$\eqalign{\theta A
&\le\sum_{m=0}\sp{\infty}\lambda I_{k_m,l_m}\cr
&\le\sum_{n=0}\sp{\infty}\sum_{j=0}\sp{\infty}\lambda I_{nj}\cr
&\le\sum_{n=0}\sp{\infty}(\theta A_n+2^{-n}\epsilon)\cr
&=\sum_{n=0}\sp{\infty}\theta A_n+\sum_{n=0}\sp{\infty}2^{-n}\epsilon\cr
&=\sum_{n=0}\sp{\infty}\theta A_n+2\epsilon.\cr}$$

\noindent Karena $\epsilon$ sebarang,
$\theta A\le\sum_{n=0}\sp{\infty}\theta A_n$ (sekali lagi, Anda
sebaiknya memeriksa bahwa ini sah, baik
$\sum_{n=0}\sp{\infty}\theta A_n$ berhingga maupun tidak).
Karena $\langle A_n\rangle_{n\in\Bbb N}$ sebarang, $\theta$ adalah
ukuran luar.

\medskip

{\bf (b)} Karena kita selalu dapat mengambil $I_0=I$ dan
$I_j=\emptyset$ untuk $j\ge 1$, untuk memperoleh barisan interval
setengah terbuka yang menutupi $I$ dengan
$\sum_{j=0}\sp{\infty}\lambda I_j=\lambda I$, kita tentu mempunyai
$\theta I\le\lambda I$. Untuk ketaksamaan sebaliknya, gunakan 114B:
jika $I\subseteq\bigcup_{j\in\Bbb N}I_j$, maka
$\lambda I\le\sum_{j=0}^{\infty}\lambda I_j$; karena
$\langle I_j\rangle_{j\in\Bbb N}$ sebarang,
$\theta I\ge\lambda I$ dan $\theta I=\lambda I$, sebagaimana
diperlukan.
}%end of proof of 114D

\cmmnt{\medskip

\noindent{\bf Catatan} Terdapat suatu peralihan yang kurang luwes dalam
argumen (a-iv) di atas, pada tahap

\Centerline{`$\theta A\le\sum_{m=0}^{\infty}\lambda
I_{k_m,l_m}\le\sum_{n=0}^{\infty}\sum_{j=0}^{\infty}\lambda I_{nj}$'.}

\noindent Saya kira Anda akan memercayai saya seandainya saya
menghilangkan $k_m$, $l_m$ sama sekali dan sekadar menulis
`karena $A\subseteq\bigcup_{n,j\in\Bbb N}I_{nj}$, maka
$\theta A\le\sum_{n=0}^{\infty}\sum_{j=0}^{\infty}\lambda I_{nj}$'.
Saya harap Anda tidak terlalu berkecil hati jika saya mengatakan bahwa
lompatan semacam itu tidak sepenuhnya aman. Alasan saya menyisipkan nama
untuk suatu bijeksi antara $\Bbb N$ dan $\Bbb N\times\Bbb N$, serta
mengambil beberapa baris untuk menyatakan secara eksplisit bahwa
$\sum_{m=0}^{\infty}\lambda
I_{k_m,l_m}\le\sum_{n=0}^{\infty}\sum_{j=0}^{\infty}\lambda I_{nj}$,
adalah sebagai berikut. Pertama-tama, ada persoalan formal bahwa
definisi 114C menuntut suatu barisan tunggal, bukan barisan ganda.
Benarkah jelas bahwa hal itu tidak berpengaruh di sini? Jika demikian,
mengapa? Tidak mungkin ada cara untuk membenarkan peralihan itu tanpa
bergantung pada fakta bahwa $\Bbb N\times\Bbb N$ terhitung dan setiap
$\lambda I_{nj}$ tak-negatif.
Jika salah satu fakta tersebut tidak benar, metode ini akan sangat
berisiko gagal.

Pada suatu saat kita tentu perlu membahas jumlah atas himpunan indeks
tak hingga selain $\Bbb N$, termasuk himpunan indeks tak terhitung. Saya
telah menyinggungnya dalam 112Bd, dan akan kembali membahasnya dalam
226A di Jilid 2.
Untuk saat ini, saya merasa kita sudah mempunyai cukup banyak gagasan
baru untuk ditangani, dan yang kita perlukan di sini adalah siasat yang
cukup jujur untuk menangani persoalan yang langsung berada di hadapan
kita.

Anda mungkin telah memperhatikan, atau menduga, bahwa beberapa
ketaksamaan `$\le$' di sini sebenarnya harus merupakan kesamaan;
jika demikian, periksa dugaan Anda dalam 114Ya.
}%end of comment

\leader{114E}{Definisi} Karena ukuran luar Lebesgue\cmmnt{ (114C)}
adalah\cmmnt{ memang} ukuran luar\cmmnt{ (114Da)}, kita dapat
menggunakannya untuk mengonstruksi suatu ukuran $\mu$, dengan metode
\Caratheodory\cmmnt{ (113C)}.
Ukuran ini adalah {\bf ukuran Lebesgue pada $\Bbb R$}.
Himpunan-himpunan $E$ yang diukur oleh $\mu$\cmmnt{ (yakni, yang
memenuhi $\theta(A\cap E)+\theta(A\setminus E)=\theta A$ untuk setiap
$A\subseteq\Bbb R$)} disebut {\bf terukur Lebesgue}.

Himpunan-himpunan yang terabaikan untuk $\mu$ disebut {\bf terabaikan
Lebesgue}\cmmnt{; perhatikan bahwa himpunan-himpunan ini tepatlah
himpunan $A$ yang memenuhi $\theta A=0$, dan semuanya terukur Lebesgue
(113Xa)}.



\leader{114F}{Lema} Misalkan $x\in\Bbb R$. Maka
$H_x=\ooint{-\infty,x}$ terukur Lebesgue untuk setiap $x\in\Bbb R$.

\proof{{\bf (a)} Intinya adalah bahwa
$\lambda I=\lambda(I\cap H_x)+\lambda(I\setminus H_x)$ untuk setiap
interval setengah terbuka $I\subseteq\Bbb R$. \Prf\ Jika
$I\subseteq H_x$ atau $I\cap H_x=\emptyset$, hal ini langsung jelas.
Jika tidak, $I$ harus berbentuk $\coint{a,b}$, dengan $a<x<b$.
Sekarang $I\cap H_x=\coint{a,x}$ dan
$I\setminus H_x=\coint{x,b}$ keduanya merupakan interval setengah
terbuka, dan

\Centerline{$\lambda(I\cap H_x)+\lambda(I\setminus H_x)=(x-a)+(b-x)
=b-a=\lambda I$. \Qed}

\medskip

{\bf (b)} Sekarang andaikan bahwa $A$ adalah sembarang subhimpunan
$\Bbb R$, dan $\epsilon>0$. Maka kita dapat menemukan barisan
$\langle I_j\rangle_{j\in\Bbb N}$ yang terdiri atas interval setengah
terbuka sedemikian sehingga
$A\subseteq\bigcup_{j\in\Bbb N}I_j$ dan
$\sum_{j=0}^{\infty}\lambda I_j\le\theta A+\epsilon$. Sekarang
$\langle I_j\cap H_x\rangle_{j\in\Bbb N}$ dan
$\langle I_j\setminus H_x\rangle_{j\in\Bbb N}$ adalah barisan
interval setengah terbuka, serta
$A\cap H_x\subseteq\bigcup_{j\in\Bbb N}(I_j\cap H_x)$,
$A\setminus H_x\subseteq\bigcup_{j\in\Bbb N}(I_j\setminus H_x)$.
Jadi

$$\eqalign{\theta(A\cap H_x)+\theta(A\setminus H_x)
&\le\sum_{j=0}^{\infty}\lambda(I_j\cap
H_x)+\sum_{j=0}^{\infty}\lambda(I_j\setminus H_x) \cr
&=\sum_{j=0}^{\infty}\lambda I_j
\le\theta A+\epsilon.\cr}$$

\noindent Karena $\epsilon$ sebarang,
$\theta(A\cap H_x)+\theta(A\setminus H_x)\le\theta A$;
karena $A$ sebarang, $H_x$ terukur, sebagaimana dicatat dalam 113D.
}%end of proof of 114F

\leader{114G}{Proposisi} Semua subhimpunan Borel dari $\Bbb R$ terukur
Lebesgue; khususnya, semua himpunan terbuka dan semua himpunan dari
kelas-kelas berikut, beserta gabungan terhitung dari himpunan-himpunan
tersebut:

(i) interval terbuka $\left]a,b\right[$,
$\left]-\infty,b\right[$, $\left]a,\infty\right[$,
$\left]-\infty,\infty\right[$, dengan $a<b\in\Bbb R$;

(ii) interval tertutup $[a,b]$, dengan $a\le b\in\Bbb R$;

(iii) interval setengah terbuka $\left[a,b\right[$,
$\left]a,b\right]$, $\left]-\infty,b\right]$,
$\left[a,\infty\right[$, dengan $a<b$ dalam $\Bbb R$.

\noindent Kita\cmmnt{ juga} mempunyai rumus berikut untuk ukuran
himpunan-himpunan tersebut, dengan $\mu$ menyatakan ukuran Lebesgue:

\Centerline{$\mu\left]a,b\right[=\mu[a,b]
=\mu\left[a,b\right[=\mu\left]a,b\right]=b-a$}

\noindent setiap kali $a\le b$ dalam $\Bbb R$, sedangkan semua interval
tak berbatas mempunyai ukuran tak hingga. Akibatnya, setiap subhimpunan
terhitung dari $\Bbb R$ terukur dan berukuran nol.


\proof{{\bf (a)} Pertama-tama saya tunjukkan bahwa semua subhimpunan
terbuka dari $\Bbb R$ terukur. \Prf\ Misalkan
$G\subseteq\Bbb R$ terbuka. Misalkan $K\subseteq\Bbb Q\times\Bbb Q$
adalah himpunan pasangan $(q,q')$ bilangan rasional sedemikian sehingga
$\coint{q,q'}\subseteq G$. Sekarang, menurut catatan dalam 111E-111F
--- khususnya, 111Eb, yang menunjukkan bahwa $\Bbb Q$ terhitung,
111F(b-iii), yang menunjukkan bahwa produk himpunan-himpunan terhitung
adalah terhitung, dan 111F(b-i), yang menunjukkan bahwa subhimpunan dari
himpunan terhitung adalah terhitung --- kita melihat bahwa $K$
terhitung. Selain itu, setiap $\coint{q,q'}$ terukur, karena merupakan
$H_{q'}\setminus H_q$ dalam notasi 114F. Jadi, menurut 111Fa,
$G'=\bigcup_{(q,q')\in K}\coint{q,q'}$ terukur.

Menurut definisi $K$, $G'\subseteq G$. Di pihak lain, jika $x\in G$,
terdapat $\epsilon>0$ sedemikian sehingga
$\ooint{x-\epsilon,x+\epsilon}\subseteq G$.
Sekarang terdapat bilangan-bilangan rasional
$q\in\ocint{x-\epsilon,x}$ dan $q'\in\ocint{x,x+\epsilon}$, sehingga
$(q,q')\in K$ dan $x\in\coint{q,q'}\subseteq G'$.
Karena $x$ sebarang, $G=G'$ dan $G$ terukur.\ \Qed

\medskip

{\bf (b)} Sekarang keluarga $\Sigma$ dari himpunan-himpunan terukur
Lebesgue adalah aljabar-$\sigma$ atas subhimpunan-subhimpunan $\Bbb R$
yang memuat keluarga himpunan terbuka, sehingga harus memuat setiap
himpunan Borel, menurut definisi himpunan Borel (111G).
\medskip

{\bf (c)} Di antara jenis-jenis interval yang ditinjau, semua interval
terbuka memang merupakan himpunan terbuka, sehingga tentu merupakan
himpunan Borel. Komplemen suatu interval tertutup dapat dinyatakan
sebagai gabungan paling banyak dua interval terbuka, sehingga merupakan
himpunan Borel, dan interval tertutup itu sendiri, sebagai komplemen
suatu himpunan Borel, juga merupakan himpunan Borel. Interval setengah
terbuka yang berbatas dapat dinyatakan sebagai irisan suatu interval
terbuka dengan suatu interval tertutup, sehingga merupakan himpunan
Borel; akhirnya, interval tak berbatas yang berbentuk
$\left]-\infty,b\right]$ atau $\left[a,\infty\right[$ adalah komplemen
suatu interval terbuka, sehingga juga merupakan himpunan Borel.

\medskip

{\bf (d)} Untuk menghitung ukurannya, dari 114Db kita sudah mengetahui
bahwa

\Centerline{$\mu\coint{a,b}=\theta\coint{a,b}=b-a$}

\noindent jika $a\le b$. Untuk jenis-jenis interval berbatas lainnya,
cukup diperhatikan bahwa $\mu\{a\}=0$ untuk setiap $a\in\Bbb R$, sebab
interval-interval tersebut hanya berbeda pada satu atau dua titik; dan
ini benar karena $\{a\}\subseteq\coint{a,a+\epsilon}$,
sehingga $\mu\{a\}\le\epsilon$, untuk setiap $\epsilon>0$.

Adapun interval-interval tak berbatas, interval-interval itu memuat
interval setengah terbuka yang panjangnya sebesar apa pun, sehingga
harus mempunyai ukuran tak hingga.

\medskip

{\bf (e)} Sebagaimana baru saja dicatat, $\mu\{a\}=0$ untuk setiap
$a$. Jika $A\subseteq\Bbb R$ terhitung, maka himpunan tersebut kosong
atau dapat dinyatakan sebagai $\{a_n:n\in\Bbb N\}$. Dalam kasus pertama,
$\mu A=\mu\emptyset=0$; dalam kasus kedua,
$A=\bigcup_{n\in\Bbb N}\{a_n\}$ adalah himpunan Borel dan
$\mu A\le\sum_{n=0}^{\infty}\mu\{a_n\}=0$.
}%end of proof of 114G

\exercises{
\leader{114X}{Latihan dasar $\pmb{>}$(a)} Misalkan
$g:\Bbb R\to\Bbb R$ sembarang fungsi tak-menurun. Untuk interval
setengah terbuka $I\subseteq\Bbb R$, definisikan $\lambda_gI$ dengan
menetapkan

\Centerline{$\lambda_g\emptyset=0$,\quad
$\lambda_g\coint{a,b}
=\lim_{x\uparrow b}g(x)-\lim_{x\uparrow a}g(x)$}

\noindent jika $a<b$. Untuk setiap himpunan $A\subseteq\Bbb R$,
tetapkan

$$\eqalign{\theta_gA=\inf\{\sum_{j=0}\sp{\infty}\lambda_gI_j:\langle
I_j\rangle_{j\in\Bbb N}
\text{ adalah barisan }&\text{interval setengah terbuka}\cr
&\text{sedemikian sehingga }
A\subseteq\bigcup_{j\in\Bbb N}I_j\}.\cr}$$

\noindent Tunjukkan bahwa $\theta_g$ adalah ukuran luar pada $\Bbb R$.
Misalkan $\mu_g$ ukuran yang didefinisikan dari $\theta_g$ dengan metode
\Caratheodory; tunjukkan bahwa $\mu_gI$ terdefinisi dan sama dengan
$\lambda_gI$ untuk setiap interval setengah terbuka
$I\subseteq\Bbb R$, dan bahwa setiap subhimpunan Borel dari $\Bbb R$
termasuk dalam domain $\mu_g$.

($\mu_g$ adalah ukuran {\bf Lebesgue-Stieltjes} yang terkait dengan
$g$.)

\spheader 114Xb Pada titik mana argumen 114Xa akan gagal jika kita
menuliskan $\lambda_g\coint{a,b}=g(b)-g(a)$ alih-alih menggunakan rumus
yang diberikan?

\sqheader 114Xc Tuliskan $\theta$ untuk ukuran luar Lebesgue dan $\mu$
untuk ukuran Lebesgue pada $\Bbb R$. Tunjukkan bahwa
$\theta A=\inf\{\mu E:E$ terukur Lebesgue, $A\subseteq E\}$ untuk
setiap $A\subseteq\Bbb R$.
\Hint{Tinjau himpunan-himpunan $E$ berbentuk
$\bigcup_{j\in\Bbb N}I_j$, dengan $\sequence{j}{I_j}$ suatu barisan
interval setengah terbuka.}

\spheader 114Xd Misalkan $X$ suatu himpunan, $\Cal I$ suatu keluarga
subhimpunan $X$ sedemikian sehingga $\emptyset\in\Cal I$, dan
$\lambda:\Cal I\to\coint{0,\infty}$ suatu fungsi sedemikian sehingga
$\lambda\emptyset=0$. Definisikan $\theta:\Cal PX\to[0,\infty]$ dengan
menetapkan

\Centerline{$\theta A=\inf\{\sum_{j=0}\sp{\infty}\lambda I_j:\langle
I_j\rangle_{j\in\Bbb N}
\text{ adalah barisan di dalam }\Cal I\text{ sedemikian sehingga }
A\subseteq\bigcup_{j\in\Bbb N}I_j\}$,}

\noindent dengan menafsirkan $\inf\emptyset$ sebagai $\infty$, sehingga
$\theta A=\infty$ jika $A$ tidak ditutupi oleh barisan mana pun di
dalam $\Cal I$. Tunjukkan bahwa $\theta$ adalah ukuran luar pada $X$.

\spheader 114Xe Misalkan $E\subseteq\Bbb R$ suatu himpunan yang
mempunyai ukuran berhingga menurut ukuran Lebesgue $\mu$. Tunjukkan
bahwa untuk setiap
$\epsilon>0$ terdapat keluarga saling lepas $I_0,\ldots,I_n$ yang
terdiri atas interval setengah terbuka sedemikian sehingga
$\mu(E\symmdiff\bigcup_{j\le n}I_j)\le\epsilon$. \Hint{misalkan
$\langle J_j\rangle_{j\in\Bbb N}$ suatu barisan interval setengah
terbuka sedemikian sehingga
$E\subseteq\bigcup_{j\in\Bbb N}J_j$ dan
$\sum_{j=0}^{\infty}\mu J_j\le\mu E+\bover{1}{2}\epsilon$. Sekarang
ambil $m$ yang cukup besar dan nyatakan $\bigcup_{j\le m}J_j$ sebagai
gabungan saling lepas interval-interval setengah terbuka.}

\sqheader 114Xf Tuliskan $\theta$ untuk ukuran luar Lebesgue dan $\mu$
untuk ukuran Lebesgue pada $\Bbb R$. Andaikan $c\in\Bbb R$. Tunjukkan
bahwa $\theta(A+c)=\theta A$ untuk setiap $A\subseteq\Bbb R$, dengan
$A+c=\{x+c:x\in A\}$. Tunjukkan bahwa jika
$E\subseteq \Bbb R$ terukur, maka $E+c$ juga terukur, dan dalam hal ini
$\mu(E+c)=\mu E$.

\spheader 114Xg Tuliskan $\theta$ untuk ukuran luar Lebesgue dan $\mu$
untuk ukuran Lebesgue pada $\Bbb R$. Andaikan $c>0$. Tunjukkan bahwa
$\theta(cA)=c\theta(A)$ untuk setiap $A\subseteq\Bbb R$, dengan
$cA=\{cx:x\in A\}$.
Tunjukkan bahwa jika $E\subseteq \Bbb R$ terukur, maka $cE$ juga
terukur, dan dalam hal ini $\mu(cE)=c\mu E$.

\leader{114Y}{Latihan lanjutan (a)}
%\spheader 114Ya
Dalam (a-iv) dari bukti 114D, tunjukkan bahwa
$\sum_{m=0}^{\infty}\lambda I_{k_m,l_m}$ sebenarnya sama dengan
$\sum_{n=0}^{\infty}\sum_{j=0}^{\infty}\lambda I_{nj}$.

\spheader 114Yb Misalkan $g$, $h:\Bbb R\to\Bbb R$ dua fungsi
tak-menurun, dengan jumlah $g+h$; misalkan
$\mu_g$, $\mu_h$, $\mu_{g+h}$ ukuran-ukuran Lebesgue-Stieltjes yang
bersesuaian (114Xa). Tunjukkan bahwa

\Centerline{$\dom\mu_{g+h}=\dom\mu_g\cap\dom\mu_h$,
\quad$\mu_{g+h}E=\mu_gE+\mu_hE$ untuk setiap $E\in\dom\mu_{g+h}$.}
\spheader 114Yc Misalkan $\sequencen{g_n}$ suatu barisan fungsi
tak-menurun dari $\Bbb R$ ke $\Bbb R$, dan andaikan bahwa
$g(x)=\sum_{n=0}^{\infty}g_n(x)$ terdefinisi dan berhingga untuk setiap
$x\in\Bbb R$. Misalkan $\mu_{g_n}$, $\mu_g$ ukuran-ukuran
Lebesgue-Stieltjes yang bersesuaian. Tunjukkan bahwa

\Centerline{$\dom\mu_g=\bigcap_{n\in\Bbb N}\dom\mu_{g_n}$,
\quad$\mu_gE=\sum_{n=0}^{\infty}\mu_{g_n}E$ untuk setiap $E\in\dom\mu_g$.}

\spheader 114Yd(i) Tunjukkan bahwa jika $A\subseteq\Bbb R$ dan
$\epsilon>0$, terdapat himpunan terbuka $G\supseteq A$ sedemikian
sehingga $\theta G\le\theta A+\epsilon$, dengan $\theta$ ukuran luar
Lebesgue. (ii) Tunjukkan bahwa jika $E\subseteq\Bbb R$ terukur Lebesgue
dan $\epsilon>0$, terdapat himpunan terbuka $G\supseteq E$ sedemikian
sehingga $\mu(G\setminus E)\le\epsilon$, dengan $\mu$ ukuran Lebesgue.
\Hint{tinjau terlebih dahulu kasus $E$ berbatas.} (iii) Tunjukkan bahwa
jika $E\subseteq\Bbb R$ terukur Lebesgue, terdapat himpunan-himpunan
Borel $H_1$, $H_2$ sedemikian sehingga $H_1\subseteq E\subseteq H_2$
dan $\mu(H_2\setminus E)=\mu(E\setminus H_1)=0$.
\Hint{gunakan (ii) untuk menemukan $H_2$, lalu tinjau komplemen $E$.}

\spheader 114Ye Tuliskan $\theta$ untuk ukuran luar Lebesgue pada
$\Bbb R$. Tunjukkan bahwa suatu himpunan $E\subseteq\Bbb R$ terukur
Lebesgue jika dan hanya jika
$\theta([-n,n]\cap E)+\theta([-n,n]\setminus E)=2n$ untuk setiap
$n\in\Bbb N$. \Hint{Gunakan 114Yd untuk menunjukkan bahwa bagi setiap
$n$ terdapat himpunan-himpunan terukur $F_n$, $H_n$ sedemikian sehingga
$F_n\subseteq[-n,n]\cap E\subseteq H_n$ dan $H_n\setminus F_n$
terabaikan.}

\spheader 114Yf Ulangi 114Xc dan 114Yd-114Ye untuk ukuran-ukuran
Lebesgue-Stieltjes dari 114Xa.

\spheader 114Yg Tuliskan $\Cal B$ untuk aljabar-$\sigma$ yang terdiri
atas subhimpunan-subhimpunan Borel dari $\Bbb R$, dan misalkan
$\nu:\Cal B\to[0,\infty]$ suatu ukuran. Misalkan
$g$, $\lambda_g$, $\theta_g$, dan $\mu_g$ seperti dalam 114Xa.
Tunjukkan bahwa jika $\nu I=\lambda_gI$ untuk setiap interval setengah
terbuka $I$, maka $\nu E=\mu_gE$ untuk setiap $E\in\Cal B$.
\Hint{tinjau terlebih dahulu himpunan-himpunan terbuka $E$, kemudian
gunakan 114Yd(i) sebagaimana diperluas dalam 114Yf.}

\spheader 114Yh Tuliskan $\Cal B$ untuk aljabar-$\sigma$ yang terdiri
atas subhimpunan-subhimpunan Borel dari $\Bbb R$, dan misalkan
$\nu:\Cal B\to[0,\infty]$ suatu ukuran sedemikian sehingga
$\nu[-n,n]<\infty$ untuk setiap $n\in\Bbb N$. Tunjukkan bahwa terdapat
fungsi $g:\Bbb R\to\Bbb R$ yang tak-menurun, kontinu dari kiri, dan
sedemikian sehingga $\nu E=\mu_gE$ untuk setiap $E\in\Cal B$, dengan
$\mu_g$ didefinisikan seperti dalam 114Xa. Apakah $g$ unik?

\spheader 114Yi Tuliskan $\Cal B$ untuk aljabar-$\sigma$ yang terdiri
atas subhimpunan-subhimpunan Borel dari $\Bbb R$, dan misalkan $\nu_1$, $\nu_2$
ukuran-ukuran dengan domain $\Cal B$ sedemikian sehingga
$\nu_1 I=\nu_2I<\infty$ untuk setiap interval setengah terbuka
$I\subseteq\Bbb R$. Tunjukkan bahwa $\nu_1E=\nu_2E$ untuk setiap
$E\in\Cal B$.

\spheader 114Yj Misalkan $\Cal E$ sembarang keluarga interval setengah
terbuka dalam $\Bbb R$. Tunjukkan bahwa (i) terdapat
$\Cal C\subseteq\Cal E$ yang terhitung sedemikian sehingga
$\bigcup\Cal E=\bigcup\Cal C$ (definisi:
1A1F) (ii) bahwa $\bigcup\Cal E$ adalah himpunan Borel, sehingga
terukur Lebesgue (iii) bahwa terdapat barisan saling lepas
$\sequencen{I_n}$ yang terdiri atas interval setengah terbuka dalam
$\Bbb R$ sedemikian sehingga
$\bigcup\Cal E=\bigcup_{n\in\Bbb N}I_n$.

\spheader 114Yk Tunjukkan bahwa untuk hampir setiap $x\in\Bbb R$
(menurut ukuran Lebesgue), himpunan

\Centerline{$\{(m,n):m\in\Bbb Z,\,n\in\Bbb N\setminus\{0\},
\,|x-\Bover{m}{n}|\le\Bover{1}{n^3}\}$}

\noindent berhingga. \Hint{taksir ukuran luar dari
$\bigcup_{n\ge n_0}\bigcup_{|m|\le kn}
[\bover{m}{n}-\bover1{n^3},\bover{m}{n}+\bover1{n^3}]$ untuk $n_0$,
$k\ge 1$.} Ulangi dengan $2+\epsilon$ menggantikan $3$.

\spheader 114Yl Tuliskan $\mu$ untuk ukuran Lebesgue pada $\Bbb R$.
Tunjukkan bahwa terdapat keluarga terhitung $\Cal F$ yang terdiri atas
subhimpunan-subhimpunan terukur Lebesgue dari $\Bbb R$ sedemikian
sehingga setiap kali $\mu E$ terdefinisi dan berhingga, serta
$\epsilon>0$, terdapat $F\in\Cal F$ sedemikian sehingga
$\mu(E\symmdiff F)\le\epsilon$. \Hint{dalam 114Xe, tunjukkan bahwa kita
dapat mengambil $I_j$ yang titik-titik ujungnya rasional.}
%114Xe
}%end of exercises

\endnotes{
\Notesheader{114} Minat saya sendiri terletak pada teori ukuran
\lq abstrak', dan dari sudut pandang struktur karya ini, tujuan utama
bagian ini adalah menguraikan suatu ruang ukur taktrivial untuk menjadi
titik pusat bagi teorema-teorema umum yang menyusul. Izinkan saya
merinci metode-metode untuk mengonstruksi ruang ukur yang sudah tersedia
bagi kita. (a) Kita mempunyai ukuran bertumpu pada titik dari 112Bd;
dalam beberapa hal, ukuran-ukuran ini trivial; tetapi ukuran tersebut
memang muncul dalam penerapan, dan justru karena pada umumnya mudah
ditangani, sering kali tepat untuk menguji gagasan baru apa pun pada
ukuran-ukuran tersebut. (b) Kita mempunyai ukuran Lebesgue pada
$\Bbb R$; perumuman langsung dari konstruksi ini menghasilkan
ukuran-ukuran Lebesgue-Stieltjes (114Xa). (c) Selanjutnya, kita
mempunyai cara-cara membangun ukuran baru dari ukuran lama, dimulai
dengan ukuran subruang (113Yb), ukuran bayangan (112Xf), dan jumlah
ukuran (112Yf). Barangkali yang terpenting di antaranya ialah
`ukuran Lebesgue pada $[0,1]$', yang untuk sementara saya sebut
$\mu_1$, dengan domain $\mu_1$ berupa
$\{E:E\subseteq[0,1]$ terukur Lebesgue$\}
=\{E\cap[0,1]:E\subseteq\Bbb R$ terukur Lebesgue$\}$, dan
$\mu_1E$ tidak lain adalah ukuran Lebesgue dari $E$ untuk setiap
$E\in\dom\mu_1$. Sebenarnya, ukuran-ukuran bayangan dari ukuran Lebesgue
pada $[0,1]$ mencakup proporsi yang sangat besar dari ukuran-ukuran
peluang (yakni, ukuran yang memberi ukuran $1$ kepada seluruh ruang)
yang penting dalam penerapan biasa.

Tentu saja ukuran Lebesgue bukan hanya contoh penuntun yang dominan bagi
teori ukuran umum, tetapi juga merupakan ukuran individual yang paling
penting bagi penerapan. Karena alasan ini, kita mungkin saja --- meski
menurut saya berpandangan sempit --- membaca Bab 12-13 dalam jilid ini
dan sebagian besar Jilid 2 seolah-olah semuanya hanya berlaku bagi
ukuran Lebesgue pada $\Bbb R$. Sesungguhnya, dalam konteks inilah
sebagian besar hasil tersebut pertama kali dikembangkan. Namun, saya
percaya bahwa dalam matematika sering kali pemahaman seseorang atas
suatu konstruksi tertentu diperdalam dan diperkuat oleh pengenalan
terhadap objek-objek terkait, dan bahwa salah satu jalan untuk
menghayati hakikat ukuran Lebesgue adalah melalui kajian sifat-sifatnya
dalam konteks teori ukuran umum yang lebih abstrak.

Untuk setiap penyelidikan yang layak mengenai penerapan teori ukuran
Lebesgue, kita harus menunggu hingga Jilid 2. Namun, saya menyertakan
114Yk sebagai petunjuk mengenai salah satu cara teori ini dapat
digunakan.
}%end of notes
\discrpage
